\chapter{Technical Requirements}
\label{ch:technical-requirements}

\section{Technology stack}

\begin{longtable}{p{3.4cm}p{4.4cm}p{6.6cm}}
\caption{Technology stack (declared in package manifests)}\\
\toprule
\textbf{Category} & \textbf{Technology} & \textbf{Detail} \\
\midrule

Language & TypeScript 5.5 & Strict mode, ES2020 target, CommonJS, declaration + source maps. \\
Runtime & Node.js $\geq$ 18 & Declared in \texttt{@expo-snips/cli} engines. \\
Workspace & pnpm workspaces & \texttt{pnpm-workspace.yaml} lists \texttt{packages/*}. \\
Engine & Handlebars 4.7 & Only runtime dependency of \texttt{@expo-snips/core}. \\
CLI & commander 12 & Command registration. \\
CLI prompts & prompts 2.4 & Interactive selection and confirmation. \\
CLI output & chalk 4.1 & Terminal colouring. \\
File I/O & fs-extra 11.2 & \texttt{ensureDir}, \texttt{readJSON}, \texttt{writeJSON}, \texttt{pathExists}. \\
Extension bundler & esbuild 0.21 & CJS bundle, \texttt{node18} target, \texttt{vscode} external. \\
Extension packaging & @vscode/vsce 2.24 & Produces the \texttt{.vsix}. \\
Validation & Node built-ins & Hand-written Draft-07 subset validator in \texttt{scripts/validate.mjs}. \\
Registry serving (dev) & Node \texttt{http} & \texttt{scripts/serve.mjs}. \\
\bottomrule
\end{longtable}

\section{Frameworks and platforms}

\begin{itemize}
  \item \textbf{Target framework:} React Native, generated through Expo. Every
        recipe declares \texttt{framework: "react-native"} and
        \texttt{runtime: "expo"} (enforced by the JSON Schema enums).
  \item \textbf{Routing:} Expo Router is the primary target
        (\texttt{app/} directory conventions); react-navigation is available for
        two recipes.
  \item \textbf{Styling:} NativeWind (Tailwind-style utility classes) or React
        Native \texttt{StyleSheet}.
  \item \textbf{Editor platform:} VS Code Extension API.
\end{itemize}

\section{Runtime requirements and entry points}

\begin{longtable}{p{4cm}p{1.9cm}p{8.3cm}}
\toprule
\textbf{Entry point} & \textbf{Runtime} & \textbf{Detail} \\
\midrule
\endhead
\texttt{packages/cli/src/index.ts} & Node & CLI binary, shebang \texttt{\#!/usr/bin/env node},
declares five subcommands. \\
\texttt{packages/vscode-extension/src/extension.ts} & VS Code & \texttt{activate()} registers commands, completion provider and status bar. \\
\texttt{packages/core/src/index.ts} & Node/VS Code & Re-exports types, context and engine. \\
\texttt{packages/registry/src/index.ts} & Node & Filesystem recipe loading; also exposes an HTTP fetch helper. \\
\texttt{scripts/build-registry.mjs} & Node & Builds the deployable registry folder. \\
\texttt{scripts/serve.mjs} & Node & Local static registry server on \texttt{PORT} (default 3000). \\
\bottomrule
\end{longtable}

\section{Architecture requirements}

\begin{itemize}
  \item A single shared generation engine must be used by all clients. The CLI and
        the extension both import \texttt{@expo-snips/core} and render with the
        same \texttt{getFilesForRecipe} function.
  \item The catalog must be consumable both as a bundle embedded in the clients and
        as static JSON over HTTP.
  \item Recipe definitions must not require filesystem access at render time; all
        \texttt{.hbs} templates are inlined into the emitted recipe JSON at build
        time (\texttt{scripts/lib/build-recipes.mjs}).
  \item There is no server-side state; all state is local to the user's project.
\end{itemize}

\section{Interfaces and APIs}

\begin{longtable}{p{3.4cm}p{11cm}}
\caption{Interface inventory}\\
\toprule
\textbf{Interface} & \textbf{Definition} \\
\midrule

CLI & Five subcommands with documented flags; see Chapter~\ref{ch:cli-reference}. \\
Registry HTTP & \texttt{GET \{base\}/index.json};
\texttt{GET \{base\}/recipes/\{url-encoded-id\}.json}. See Chapter~\ref{ch:api-documentation}. \\
Programmatic & \texttt{@expo-snips/core} exports
\texttt{renderTemplate}, \texttt{getFilesForRecipe},
\texttt{GenerationContext}, \texttt{ScreenRecipe}, \texttt{FileTemplate},
\texttt{ScreenVariant} and the option union types. \\
Extension commands & \texttt{expo-snips.addScreen} and
\texttt{expo-snips.generateFromCompletion} (the latter is not declared in
\texttt{contributes.commands}, only registered at runtime). \\
Configuration & Six VS Code settings; two environment variables; two secrets used
in CI. See Chapter~\ref{ch:configuration}. \\
\bottomrule
\end{longtable}

\section{Data requirements}

\begin{itemize}
  \item \textbf{No database.} NOT APPLICABLE.
  \item \textbf{Persistent state:} a JSON manifest (\texttt{.expo-snips.json})
        written to the user's project root. Its shape is defined in
        Section~5.8 of Chapter~\ref{ch:functional-requirements}.
  \item \textbf{Catalog data:} JSON recipe documents plus inlined Handlebars
        template strings.
  \item \textbf{Schema:} \texttt{packages/registry/schema.json}
        (JSON Schema Draft-07).
\end{itemize}

\section{Authentication and authorization requirements}

The product itself requires no authentication and implements no authorization.
\textbf{NOT APPLICABLE.} See Chapters~\ref{ch:authentication} and
\ref{ch:authorization} for why, and for the authentication templates that are
generated \emph{for} user applications.

\section{Infrastructure requirements}

\begin{itemize}
  \item \textbf{Build:} Node 18+ and pnpm 9+.
  \item \textbf{Distribution:} npm registry for the two public packages; GitHub
        Pages for the registry; manual \texttt{.vsix} for the extension.
  \item \textbf{CI:} GitHub Actions, \texttt{ubuntu-latest} runners.
  \item \textbf{No containerisation, orchestration or cloud infrastructure is
        required or present.}
\end{itemize}

\section{Storage requirements}

None beyond local disk. NOT APPLICABLE for object storage, databases or caches
(a 5-minute in-memory cache in the extension aside).

\section{Monitoring requirements}

\textbf{UNKNOWN / not implemented.} There is no monitoring, alerting, error
reporting or health endpoint.

\section{Testing requirements}

\textbf{Not implemented.} There are no test files, no test runner and no test
script. The only automated check in the repository is the recipe validator
(\texttt{scripts/validate.mjs}), which verifies recipe structure against the JSON
Schema and is exposed as \texttt{pnpm validate}. See Chapter~\ref{ch:testing}.

\section{Third-party integration requirements}

\begin{longtable}{p{4.4cm}p{5.2cm}p{4.8cm}}
\caption{Integrations}\\
\toprule
\textbf{Integration} & \textbf{Direction} & \textbf{Used for} \\
\midrule

npm registry & outbound (publish) & Publishing \texttt{@expo-snips/core} and \texttt{@expo-snips/cli}. \\
GitHub Pages & outbound (deploy) & Hosting the static registry. \\
GitHub Actions & build & CI/CD. \\
Registry HTTP endpoint & inbound (fetch) & Optional remote recipe source. \\
create-expo-app & outbound (exec) & \texttt{init} project bootstrap. \\
Package managers & outbound (exec) & Installing recipe dependencies. \\
VS Code API & host & Extension surface. \\
\bottomrule
\end{longtable}

Detailed in Chapter~\ref{ch:integrations}.
